\section{Audit of the supplied finite-exponent manuscript}
\label{sec:fable-audit}

We checked the supplied manuscript by Claude Fable 5, \emph{The circuit
prior across the whole strong prior regime: cloning and subsampling at
every $\gamma>1$}. We found no error invalidating its central finite-$\gamma$ proof. In
particular, the Gaussian change of measure, robustness of the cloned
witness, singular-vector frame net, orthonormalization of the net points,
and spectral truncation preserve the stated margin and cost bounds.
The conversion from a low-cost representative to a small bounded-weight
network uses the compression theorem proved in the earlier paper.

The main improvement concerns the learning theorem. The supplied proof
bounds the tail by multiplying the largest excluded function probability
by $2^{2^d}$, the number of Boolean functions. This gives a valid but
unnecessary additional requirement $N^\gamma\gtrsim2^d$. Its own
structural estimate yields a stronger conclusion: outside one exceptional
parameter event, \emph{every} globally margin-separated network has a
low-cost representative. The total probability of functions without such
a representative is therefore bounded by the exceptional parameter
probability directly. Counting only the resulting low-cost class gives
the width-uniform learning theorem in the accompanying paper. The source's
exponential width requirement is thus a limitation of its tail estimate,
not a demonstrated obstruction to learning.

We also replace the singular-vector net with Gaussian enlargement
applied to distance from an operator-norm ball. This shortens the proof
and removes its rank and net-entropy bookkeeping. A direct chi-square
estimate removes the logarithmic width proviso in its small-width
complexity comparison.

Two qualifications concern the source's discussion rather than its main
proof. First, the simulation in which a spectral cut lies below the noise
bulk demonstrates failure of that particular cut. It does not establish
necessity of the theorem's quantitative width threshold, nor a lower
bound on the width required for generalization. Second, the conjectured
$\gamma=1$ extension is not proved by the source's vanishing-operator-noise
argument. Any endpoint result needs a separate argument; the endpoint
statements in the accompanying endpoint paper are proved independently.

Minor points do not affect the results: ceiling operations in the
small-width bound require an additive one or an absolute factor; the
rank estimate is $r_*\le18K/\delta^2$; and the no-inner-layer case $L=2$
is most clearly handled as an empty truncation sum. We found no major
mistake invalidating the supplied manuscript's finite-$\gamma>1$
complexity theorem.
